%%%%%%%%%%%%%%%%%%%%%%%%%%%%%%%%%%%%%%%%%
% کارنامۀ کامل (نسخۀ فارسی)
%
% این نسخه، برخلاف سه نسخۀ دیگر، چیزی را حذف نمی‌کند: همۀ سوابق پژوهشی،
% کاری، آموزشی و پروژه‌ها --- از جمله مواردی که در نسخه‌های انگلیسی
% کامنت شده‌اند --- در آن آمده است.
%
% کامپایل: xelatex Kahbod-Aeini_FullCV_FA.tex
%%%%%%%%%%%%%%%%%%%%%%%%%%%%%%%%%%%%%%%%%

\documentclass[11pt]{article}
\usepackage{fa-cv}

\begin{document}

\faHeader{کهبد آیینی}{%
دانشکدۀ مهندسی برق و کامپیوتر، دانشگاه واترلو، واترلو، انتاریو، کانادا\\[0.3em]
\lr{\href{mailto:kaeini@uwaterloo.ca}{kaeini@uwaterloo.ca}} |
\lr{\href{https://www.linkedin.com/in/kahbod-aeini-833b71190/}{LinkedIn}} |
\lr{\href{https://github.com/kahbodaeini}{GitHub}} |
\lr{\href{https://kahbod.org}{kahbod.org}} |
\lr{\href{https://scholar.google.ca/citations?user=3Fd5R9YAAAAJ&hl=en}{Google Scholar}}%
}

%----------------------------------------------------------------------------------------
\begin{faSection}{دربارۀ من}
دانشجوی دکتری مهندسی برق و کامپیوتر در دانشگاه واترلو، در آزمایشگاه الگوریتم‌ها و سیستم‌های توزیع‌شده، به راهنمایی پروفسور \lr{Wojciech Golab}. پیش از آن، کارشناسی ارشد علوم کامپیوتر را در دانشگاه کلگری زیر نظر پروفسور \lr{Philipp Woelfel} و کارشناسی مهندسی کامپیوتر را در دانشگاه صنعتی شریف (آریامهر) در سال ۲۰۲۳ به پایان رساندم.

\vspace{0.4em}
کار پژوهشی من بر الگوریتم‌های حافظۀ مشترک و اولیه‌های همگام‌سازی برای سیستم‌های دارای حافظۀ پایدار (غیرفرّار) متمرکز است، با علاقۀ ویژه به اشیای بازیابی‌پذیر: ساختمان‌های دادۀ همروندی که وضعیتشان پس از خرابی سیستم سازگار می‌ماند. به طراحی الگوریتم‌های کارآمد و اثبات‌پذیر که پل میان نظریه و عمل در محاسبات توزیع‌شده باشند علاقه‌مندم.

\vspace{0.4em}
هم‌زمان با دکتری، به‌صورت پاره‌وقت مهندس نرم‌افزار شرکت \lr{ArcTrade} هستم و روی سرویس‌های سمت سرور و واسط‌های برنامه‌نویسی سامانۀ عامل خرده‌فروشی انرژی این شرکت کار می‌کنم؛ کاری که از معاملات انرژی و مدیریت ریسک، قیمت‌گذاری، صورت‌حساب و تسویه برای خرده‌فروشان برق در بازارهای آمریکای شمالی پشتیبانی می‌کند.
\end{faSection}

%----------------------------------------------------------------------------------------
\begin{faSection}{تحصیلات}

\faEntry{دانشگاه واترلو}{ژانویۀ ۲۰۲۶ تاکنون}
دکتری مهندسی برق و کامپیوتر\\
استاد راهنما: پروفسور \lr{\href{https://uwaterloo.ca/distributed-algorithms-systems-lab/}{Wojciech Golab}}\\
درس‌های مرتبط: برنامه‌نویسی چندهسته‌ای (موضوعات پیشرفتۀ پژوهشی)

\vspace{0.5em}
\faEntry{دانشگاه کلگری}{۲۰۲۳ تا ۲۰۲۵}
کارشناسی ارشد علوم کامپیوتر\\
استاد راهنما: پروفسور \lr{\href{https://www.woelfellab.org/about.html}{Philipp Woelfel}}\\
پایان‌نامه: \href{https://ucalgary.scholaris.ca/bitstreams/5ae5cf8f-5009-4db3-b91e-d3dbb679b5bf/download}{«قفل بی‌انتظارِ تصادفی با فضای بهینه»} --- طراحی و تحلیل نخستین الگوریتم قفلِ بی‌انتظارِ تصادفی با فضای بهینه برای سیستم‌های حافظۀ مشترک\\
درس‌های مرتبط: الگوریتم‌های تصادفی، الگوریتم‌های محاسبات توزیع‌شده، داده‌کاوی و یادگیری ماشین

\vspace{0.5em}
\faEntry{دانشگاه صنعتی شریف (آریامهر)}{۲۰۱۹ تا ۲۰۲۳}
کارشناسی مهندسی کامپیوتر\\
پایان‌نامه: \href{https://github.com/kahbodaeini/Fraud-Detection}{«تشخیص تقلب از روی دنباله‌های تراکنش‌های بانکی با الگوریتم‌های گوناگون یادگیری ماشین»}\\
رتبۀ ۹۴ کشوری (۳۱ منطقه) در کنکور سراسری میان حدود ۱۶۰٬۰۰۰ شرکت‌کننده\\
درس‌های مرتبط: طراحی الگوریتم‌ها، سیستم‌های پایگاه داده، زیست‌اطلاعاتی، جبر خطی، هوش مصنوعی، یادگیری ماشین، محاسبات عددی

\vspace{0.5em}
\faEntry{سازمان ملی پرورش استعدادهای درخشان (سمپاد)}{۲۰۱۳ تا ۲۰۱۹}
دیپلم دبیرستان در رشتۀ ریاضی و فیزیک --- گزینشی‌ترین برنامۀ آموزشی ایران

\end{faSection}

%----------------------------------------------------------------------------------------
\begin{faSection}{علایق پژوهشی}
اشیای بازیابی‌پذیر و الگوریتم‌های حافظۀ پایدار، ساختمان‌های دادۀ همروند و حافظۀ مشترک،
الگوریتم‌های تصادفی و نظریۀ احتمال، سیستم‌های توزیع‌شده، ذخیره‌سازهای توزیع‌شده و پایگاه‌های داده،
برنامه‌نویسی چندهسته‌ای، نظریۀ پیچیدگی، ساختمان‌های دادۀ پیشرفته
\end{faSection}

%----------------------------------------------------------------------------------------
\begin{faSection}{انتشارات}

\textbf{کهبد آیینی}، \lr{Dante Bencivenga}، \lr{George Giakkoupis}، \lr{Philipp Woelfel}\\
\href{https://dl.acm.org/doi/10.1145/3796701.3815952}{\textit{«قفل‌های بی‌انتظارِ تصادفیِ ساده و کارآمد»}}\\
\href{https://www.podc.org/podc2026/}{چهل‌وپنجمین سمپوزیوم \lr{ACM} در مبانی محاسبات توزیع‌شده (\lr{PODC})}، لندن، ژوئیۀ ۲۰۲۶. \textbf{جایزۀ بهترین مقاله.}

\vspace{0.5em}
\textbf{کهبد آیینی}، \lr{Wojciech Golab}\\
\href{https://dl.acm.org/doi/10.1145/3820355.3820388}{\textit{«تحلیل خطی‌پذیری در سیستم‌های توزیع‌شدۀ نسبیتی»}}\\
هشتمین کارگاه \lr{ApPLIED}، هم‌زمان با \lr{PODC 2026}، لندن، ژوئیۀ ۲۰۲۶.

\vspace{0.5em}
\textbf{کهبد آیینی}\\
\textit{«قفل بی‌انتظارِ تصادفی با فضای بهینه»}\\
\href{https://ucalgary.scholaris.ca/bitstreams/5ae5cf8f-5009-4db3-b91e-d3dbb679b5bf/download}{پایان‌نامۀ کارشناسی ارشد}، دانشگاه کلگری، ۲۰۲۵.

\end{faSection}

\begin{faSection}{مقالات در دست داوری}
\textbf{کهبد آیینی}، \lr{Brianna Wang}، \lr{Wojciech Golab}، \lr{Neeraj Mittal}\\
\textit{«ساخت جهانیِ سریع‌ترِ اشیای خواندن-تغییر-نوشتنِ آشکارپذیر برای حافظۀ پایدار»}\\
در دست داوری، سمپوزیوم \lr{ACM SIGPLAN} در اصول و عمل برنامه‌نویسی موازی (\lr{PPoPP})، ۲۰۲۷.
\end{faSection}

%----------------------------------------------------------------------------------------
\begin{faSection}{سوابق کاری}

\faEntry{مهندس نرم‌افزار --- \lr{ArcTrade} (پاره‌وقت، هم‌زمان با دکتری)}{دسامبر ۲۰۲۵ تاکنون}
\begin{faItems}
  \item طراحی و ساخت واسط‌های \lr{GraphQL} و \lr{REST} برای سامانۀ عامل خرده‌فروشی انرژیِ \lr{ArcTrade} (\lr{Retail eOS}).
  \item مشارکت اصلی در ماژول امنیت سکو: احراز هویت، مجوزدهی و کنترل دسترسی مبتنی بر مجوز روی داده‌های حساس بازار و مالی.
  \item نوشتن و بهینه‌سازی پرس‌وجوهای گستردۀ \lr{SQL} روی مدل دادۀ رابطه‌ای سکو، به‌عنوان پشتوانۀ \lr{resolver}های واسط برنامه‌نویسی و مسیرهای گزارش‌گیری.
  \item پیاده‌سازی تبادل الکترونیکی دادۀ امن و خودکار (\lr{EDI}) با شرکای تجاری بیرونی.
  \item مشارکت در بخش‌های پشتیبانِ معاملات انرژی و مدیریت ریسک (\lr{ETRM})، قیمت‌گذاری، صورت‌حساب و تسویه برای خرده‌فروشان برق در بازارهای آمریکای شمالی.
  \item همکاری در طراحی فرایند استخدام و آموزش نیروهای تازه‌وارد تیم.
  \item همکاری با تیم‌های میان‌رشته‌ای برای طراحی نرم‌افزار مقیاس‌پذیر و قابل‌اتکا متناسب با پیچیدگی عملیاتی و مالی صنعت انرژی.
\end{faItems}

\faEntry{مهندس داده --- \lr{AI Lida}}{فوریه تا سپتامبر ۲۰۲۲}
\begin{faItems}
  \item توسعه، نگهداری و بهینه‌سازی خطوط لولۀ داده برای فرایندهای \lr{ETL} با \lr{Python} و \lr{PostgreSQL}.
  \item همکاری با واحد امنیت برای گردآوری نیازمندی‌ها و ارائۀ راهکارهای دادهایِ هم‌راستا با اهداف کسب‌وکار.
  \item مصاحبه، استخدام، آموزش و سرپرستی تیمی پنج‌نفره از کارآموزان مهندسی داده، همراه با راهنمایی و بازخورد عملکرد.
\end{faItems}

\end{faSection}

%----------------------------------------------------------------------------------------
\begin{faSection}{سوابق پژوهشی}

\faEntry{دانشجوی دکتری و دستیار پژوهشی، الگوریتم‌های توزیع‌شده --- دانشگاه واترلو}{ژانویۀ ۲۰۲۶ تاکنون}
استاد راهنما: پروفسور \lr{Wojciech Golab}، آزمایشگاه الگوریتم‌ها و سیستم‌های توزیع‌شده\\
پژوهش در حوزۀ حافظۀ پایدار، سیستم‌های توزیع‌شدۀ مبتنی بر حافظۀ مشترک و اشیای بازیابی‌پذیر، با تمرکز بر طراحی و تحلیل الگوریتم‌های کارآمد برای همگام‌سازی، اجماع و تحمل خطا؛ ترکیب تحلیل نظری با طراحی عملی الگوریتم.

\vspace{0.5em}
\faEntry{دانشجوی کارشناسی ارشد و دستیار پژوهشی، الگوریتم‌های توزیع‌شده --- دانشگاه کلگری}{ژوئن ۲۰۲۴ تا ژوئن ۲۰۲۵}
استاد راهنما: پروفسور \lr{Philipp Woelfel}\\
طراحی و تحلیل الگوریتم‌های قفلِ بی‌انتظارِ تصادفی با موازنۀ بهینۀ فضا و زمان در سیستم‌های حافظۀ مشترک؛ از جمله نخستین الگوریتم با فضای بهینه و گونه‌های مکمل با زمان بهینه. تحلیل کارایی با نظریۀ احتمال و اثبات درستی و بی‌انتظاری.

\vspace{0.5em}
\faEntry{دستیار پژوهشی زیست‌اطلاعاتی (داوطلبانه، از راه دور) --- دانشگاه جانز هاپکینز}{فوریه تا اوت ۲۰۲۳}
استاد راهنما: پروفسور \lr{Mehdi Pirooznia}\\
انتخاب ویژگی و توسعۀ مدل‌های یادگیری ماشین برای پیش‌بینی ترکیب سلول‌های ایمنی در سرطان ریه؛ شناسایی انواع سلول‌های ایمنی، تحلیل ژن‌های با بیان متفاوت، تحلیل غنی‌سازی عملکردی و دسته‌بندی زیرگونه‌های سرطان.

\vspace{0.5em}
\faEntry{دستیار پژوهشی یادگیری ماشین --- آزمایشگاه مهندسی نرم‌افزار هوشمند (\lr{ISE})، دانشگاه صنعتی شریف}{اوت تا نوامبر ۲۰۲۲}
استاد راهنما: پروفسور عباس حیدرنوری\\
توسعۀ مدل‌های یادگیری عمیق برای پیش‌بینی کد و تحلیل معنایی برنامه‌های \lr{Python}، با استخراج ویژگی پیشرفته از دفترچه‌های \lr{Jupyter} در مجموعه‌دادۀ \lr{KGTorrent}؛ در راستای سامانه‌های خودکار کمک‌برنامه‌نویسی و درک کد.

\vspace{0.5em}
\faEntry{دستیار پژوهشی یادگیری عمیق --- پژوهشگاه دانش‌های بنیادی (\lr{IPM})}{اکتبر ۲۰۲۱ تا فوریۀ ۲۰۲۲}
اساتید راهنما: دکتر فاطمه بهاری‌فرد و دکتر وحید متقد\\
توسعه و ارزیابی الگوریتم‌های خوشه‌بندی چهره با شبکه‌های کانولوشنی گرافی (\lr{GCN}) برای پیش‌بینی پیوند و دسته‌بندی زوجی؛ پیاده‌سازی خط لولۀ ارزیابی و محک‌زنی.

\end{faSection}

%----------------------------------------------------------------------------------------
\begin{faSection}{جوایز و افتخارات}

\faEntry{جایزۀ بهترین مقاله --- چهل‌وپنجمین سمپوزیوم \lr{ACM} در مبانی محاسبات توزیع‌شده (\lr{PODC})}{۲۰۲۶}
برای مقالۀ «قفل‌های بی‌انتظارِ تصادفیِ ساده و کارآمد».

\vspace{0.4em}
\faEntry{جایزۀ پژوهشی دانشکده --- دانشکدۀ علوم کامپیوتر، دانشگاه کلگری}{۲۰۲۳}
جایزۀ دستیاری پژوهشی که به بهترین دانشجویان هر دانشکده تعلق می‌گیرد تا به‌جای وظایف آموزشی بر پژوهش خود تمرکز کنند.

\vspace{0.4em}
\faEntry{میان صد نفر برتر کشور}{۲۰۱۹}
رتبۀ ۹۴ کشوری (۳۱ منطقه) در کنکور سراسری میان حدود ۱۶۰٬۰۰۰ شرکت‌کننده.

\vspace{0.4em}
\faEntry{بورسیۀ بنیاد ملی نخبگان}{۲۰۱۹}
عضویت و بورسیه در بنیاد ملی نخبگان ایران.

\end{faSection}

%----------------------------------------------------------------------------------------
\begin{faSection}{سوابق تدریس}

\faEntry{دانشگاه کلگری}{۲۰۲۳ تا ۲۰۲۶}
سرپرستی و مدیریت تیم‌های دستیار آموزشی در درس‌های کارشناسی و تحصیلات تکمیلی (در مجموع تا ۴۳۰ دانشجو).
\begin{faItems}
  \item \textbf{مبانی نظری علوم کامپیوتر ۲} --- سرپرست دستیاران آموزشی؛ مدرس: \lr{Philipp Woelfel}؛ زمستان ۲۰۲۴، ۲۰۲۵ و ۲۰۲۶. هماهنگی ۷ دستیار آموزشی و پشتیبانی بیش از ۱۵۰ دانشجو در نظریۀ اتوماتا و احتمال.
  \item \textbf{پارادایم‌های برنامه‌نویسی} --- سرپرست دستیاران آموزشی؛ مدرس: \lr{Ahmad Nasri}؛ پاییز ۲۰۲۵. مدیریت ۶ دستیار آموزشی و برگزاری کلاس‌های حل تمرین \lr{Haskell} و \lr{Prolog} برای بیش از ۷۵ دانشجو.
  \item \textbf{الگوریتم‌های تصادفی} (درس تحصیلات تکمیلی) --- مدرس: \lr{Philipp Woelfel}؛ پاییز ۲۰۲۴. ارائۀ کلاس‌های حل تمرین در ساختمان‌های دادۀ تصادفی، تکنیک‌های کران پایین و کلاس‌های پیچیدگی.
  \item \textbf{مقدمه‌ای بر علوم کامپیوتر (برنامه‌نویسی \lr{Python})} --- مدرس: \lr{Richard Zhao}؛ پاییز ۲۰۲۳. تدریس مبانی \lr{Python}، ساختمان داده و الگوریتم‌ها برای بیش از ۵۰ دانشجو.
\end{faItems}

\faEntry{دانشگاه صنعتی شریف (آریامهر)}{۲۰۲۰ تا ۲۰۲۳}
سرپرست دستیاران آموزشی و دستیار آموزشی در چند درس کارشناسی (تیم‌های ۱۰ تا ۲۵ نفره، بیش از ۲۰۰ دانشجو).
\begin{faItems}
  \item \textbf{ساختمان داده و الگوریتم‌ها} --- سرپرست دستیاران آموزشی؛ مدرس: مهدی صفرنژاد؛ زمستان ۲۰۲۲. طراحی تمرین‌ها و پروژه‌ها، سرپرستی تیمی بیش از ۲۵ نفره.
  \item \textbf{برنامه‌نویسی موبایل} --- سرپرست دستیاران آموزشی؛ مدرس: سعید ناری؛ زمستان ۲۰۲۲. طراحی و ارزیابی پروژه‌های توسعۀ \lr{Android}.
  \item \textbf{محاسبات عددی} --- سرپرست دستیاران آموزشی؛ مدرس: فاطمه بهاری‌فرد؛ پاییز ۲۰۲۲. تدوین برنامۀ پروژۀ پایانی، سرپرستی تیمی بیش از ۱۰ نفره.
\end{faItems}

\end{faSection}

%----------------------------------------------------------------------------------------
\begin{faSection}{فعالیت‌های حرفه‌ای}
\faEntry{داور بیرونی --- سمپوزیوم بین‌المللی محاسبات توزیع‌شده (\lr{DISC})}{۲۰۲۶}
\end{faSection}

%----------------------------------------------------------------------------------------
\begin{faSection}{پروژه‌ها}

\textbf{ساختمان‌های دادۀ همروند روی سخت‌افزار چندهسته‌ای.} طراحی، پیاده‌سازی و ارزیابی درخت‌های جست‌وجوی دودویی همروند مبتنی بر \lr{TLE} و جدول‌های درهم‌سازی با کاوش گسترش‌پذیر، با حافظۀ تراکنشی محدود اینتل (\lr{RTM}) روی سروری \lr{NUMA} با ۱۹۲ نخ. تحلیل رفتار تثبیت و لغو تراکنش‌ها، اشتراک کاذب، اثرهای \lr{NUMA} و راهبردهای بازپس‌گیری حافظه (\lr{EBR}) با \lr{perf}؛ بررسی گسترش موازی در مسیر جایگزین با \lr{OpenMP} و \lr{CAS}. (درس \lr{CS 798}، زمستان ۲۰۲۶)

\vspace{0.4em}
\textbf{مدل تشخیص تقلب.} پایان‌نامۀ کارشناسی: توسعه، آموزش و آزمون سه مدل یادگیری ماشین برای تشخیص تقلب روی مجموعه‌دادۀ تراکنش‌های بانکی و مقایسۀ تحلیلی روش‌ها. (\href{https://github.com/kahbodaeini/Fraud-Detection}{گیت‌هاب} --- ژوئن ۲۰۲۳)

\vspace{0.4em}
\textbf{اپلیکیشن گفت‌وگوی امن.} توسعۀ یک اپلیکیشن گفت‌وگوی امن با \lr{Python} در تیمی سه‌نفره، محافظت‌شده با روش‌های گوناگون امنیت رایانه‌ای؛ پروژۀ پایانی درس امنیت داده و شبکه. (\href{https://github.com/rerfan24/DNS-project}{گیت‌هاب} --- ژوئن ۲۰۲۳)

\vspace{0.4em}
\textbf{آبرسان هوشمند.} همکاری در تیمی چهارنفره برای طراحی و پیاده‌سازی سامانۀ مخزن آب هوشمند، با تلفیق سخت‌افزارهایی مانند \lr{Raspberry Pi}، حسگر مجاورت و حسگرهای سطح آب، همراه با وب‌اپلیکیشنی مبتنی بر \lr{Django}. (\href{https://github.com/Parsa2820/es-project}{گیت‌هاب} --- ژوئن ۲۰۲۳)

\vspace{0.4em}
\textbf{درایو ابری.} همکاری در تیمی شش‌نفره برای طراحی و پیاده‌سازی سرویس اشتراک فایل، با بهره‌گیری از \lr{Django}، \lr{Docker} و \lr{MinIO} و رعایت متدولوژی چابک. (\href{https://github.com/sad-project}{گیت‌هاب} --- دسامبر ۲۰۲۲)

\vspace{0.4em}
\textbf{سامانۀ بازیابی اطلاعات.} طراحی و پیاده‌سازی سامانۀ بازیابی برای تعاملات توییتر و مقالات سلامت، همراه با رابط گرافیکی و دسته‌بندی و خوشه‌بندی توییت‌ها و مقالات. (\href{https://github.com/IR1401-Spring-Final-Projects/SocialMediaHealth1401-19_39}{گیت‌هاب} --- ژوئن ۲۰۲۲)

\vspace{0.4em}
\textbf{پردازش زبان طبیعی.} پیاده‌سازی روش‌های پردازش زبان طبیعی روی یک مجموعه‌دادۀ گفتاری و ارائۀ نتایج و تحلیل‌ها. (\href{https://github.com/kahbodaeini/MIR_NLP}{گیت‌هاب} --- ژوئیۀ ۲۰۲۲)

\vspace{0.4em}
\textbf{تطبیق نمایه‌های کاربران.} پیاده‌سازی کامل و جمع‌بندی مقالۀ «تطبیق نمایه‌های کاربران در شبکه‌های اجتماعی مختلف بر پایۀ شناسایی چهره»؛ پروژۀ پایانی درس محاسبات عددی. (\href{https://github.com/kahbodaeini/Users_profile_matching}{گیت‌هاب} --- ژوئیۀ ۲۰۲۱)

\vspace{0.4em}
\textbf{زیست‌اطلاعاتی.} مطالعۀ رونویسی‌شناسی سرطان خون \lr{AML} برای شناسایی ژن‌های با بیان متفاوت و مسیرهای مرتبط که می‌توانند در بروز بیماری نقش داشته باشند. (\href{https://github.com/kahbodaeini/AML_Leukemia}{گیت‌هاب} --- ژانویۀ ۲۰۲۲)

\vspace{0.4em}
\textbf{تشخیص ارقام دست‌نویس.} طراحی و توسعۀ مدلی برای تشخیص ارقام دست‌نویس. (\href{https://github.com/kahbodaeini/AI}{گیت‌هاب} --- نوامبر ۲۰۲۱)

\vspace{0.4em}
\textbf{پایگاه دادۀ آمار فوتبال.} طراحی و پیاده‌سازی پایگاه دادۀ ذخیرۀ آمار فوتبال با سامانۀ مدیریت پایگاه دادۀ \lr{MySQL}.

\vspace{0.4em}
\textbf{اپلیکیشن خدمات تعمیرگاه.} توسعۀ یک اپلیکیشن \lr{Android} برای استفاده از خدمات تعمیرگاه؛ پروژۀ پایانی درس برنامه‌نویسی موبایل. (\href{https://github.com/kahbodaeini/Garage}{گیت‌هاب} --- ژوئیۀ ۲۰۲۱)

\vspace{0.4em}
\textbf{اپلیکیشن مسیریاب.} توسعۀ یک اپلیکیشن مسیریابی در \lr{Android} با قابلیت‌های کامل و مشابه \lr{Google Maps}. (\href{https://github.com/kahbodaeini/Navigation_App}{گیت‌هاب} --- فوریۀ ۲۰۲۱)

\vspace{0.4em}
\textbf{اپلیکیشن رمزارز.} توسعۀ یک اپلیکیشن \lr{Android} برای نمایش برخط آخرین قیمت و نمودار شمعی رمزارزها. (\href{https://github.com/kahbodaeini/Cryptocurrency_App}{گیت‌هاب} --- ژانویۀ ۲۰۲۱)

\vspace{0.4em}
\textbf{فروشگاه اینترنتی.} توسعۀ یک فروشگاه اینترنتی با رابط گرافیکی و معماری کارخواه-کارساز به زبان \lr{Java}. (\href{https://github.com/AdvProg2020/Project_team-39}{گیت‌هاب} --- ژوئن ۲۰۲۰)

\vspace{0.4em}
\textbf{اپلیکیشن گفت‌وگو.} توسعۀ یک اپلیکیشن گفت‌وگو به زبان \lr{C}. (\href{https://github.com/kahbodaeini/Chat_Application}{گیت‌هاب} --- ژانویۀ ۲۰۲۰)

\end{faSection}

%----------------------------------------------------------------------------------------
\begin{faSection}{درس‌های مرتبط و نمرات}
\textbf{برنامه‌نویسی چندهسته‌ای (موضوعات پیشرفتۀ پژوهشی)}: \lr{A+} --- \textbf{الگوریتم‌های تصادفی}: \lr{A+} --- \textbf{الگوریتم‌های محاسبات توزیع‌شده}: \lr{A+} --- \textbf{داده‌کاوی و یادگیری ماشین}: \lr{A} --- \textbf{بازیابی اطلاعات پیشرفته}: \lr{A+} --- \textbf{طراحی پایگاه داده}: \lr{A} --- \textbf{طراحی الگوریتم‌ها}: \lr{A} --- \textbf{هوش مصنوعی}: \lr{A+} --- \textbf{ساختمان داده و الگوریتم‌ها}: \lr{A+} --- \textbf{جبر خطی}: \lr{A+} --- \textbf{محاسبات عددی}: \lr{A+} --- \textbf{سیستم‌های بی‌درنگ}: \lr{A}
\end{faSection}

%----------------------------------------------------------------------------------------
\begin{faSection}{مهارت‌ها}

\textbf{نظریه و الگوریتم}: الگوریتم‌های توزیع‌شده، همروندی، تصادفی‌سازی، نظریۀ پیچیدگی، طراحی و بهینه‌سازی ساختمان داده

\vspace{0.3em}
\textbf{زبان‌های برنامه‌نویسی}: \lr{C\#}، \lr{Python}، \lr{SQL}، \lr{GraphQL}، \lr{R}، \lr{Haskell}، \lr{Java}، \lr{C/C++}، \lr{Prolog}

\vspace{0.3em}
\textbf{چارچوب‌ها و ابزارها}: \lr{.NET / ASP.NET Core}، \lr{Entity Framework}، \lr{Docker}، \lr{Git}، \lr{Django}، \lr{PyTorch}، \lr{TensorFlow}، \lr{scikit-learn}، \lr{NumPy}، \lr{Pandas}، \lr{OpenCV}، \lr{Android}، \lr{limma}، \lr{OpenMP}، \lr{perf}

\vspace{0.3em}
\textbf{پایگاه‌های داده}: \lr{Microsoft SQL Server}، \lr{PostgreSQL}، \lr{MongoDB}، \lr{MySQL}

\vspace{0.3em}
\textbf{سایر}: \lr{Linux}، \lr{NixOS}، \lr{Kali}

\end{faSection}

\end{document}
